% GENERATED FILE. Edit canonical lessons, then run npm run generate:books.
% canonical-source-hash: fnv1a64:11fea893a0e21f2c
% canonical-lessons: MR-C494-nate, MR-C494-natevaik, MR-C494-parivar, MR-C494-anath, MR-C494-sakkha

\chapter{Family Ties}
\label{ch:mr-family-ties}
\hlchaptermodality{494}

\begin{chapteropening}
\textbf{By the end of this chapter:} \emph{I can talk about a relationship, relatives, a household and an orphan.}

Complete the last lesson of chapter 494: I can talk about a relationship, relatives, a household and an orphan.
\end{chapteropening}

\section[\texorpdfstring{\mr{नाते} (nāte)}{नाते (nāte)}]{\mr{नाते} (nāte) — a relationship; kinship}

\label{lesson:MR-C494-nate}



\begin{quote}
Before the new one: say the Marathi for to shiver with cold, then the Marathi for to limp.
\end{quote}

\subsection*{You'll want to know: \mr{नाते}}
\textbf{\mr{नाते}} — \emph{nāte} — \textquotedblleft{}a relationship; kinship\textquotedblright{}.

\begin{quote}
\mr{आमचं} \mr{आणि} \mr{त्यांचं} \mr{जवळचं} \mr{नातं} \mr{आहे}.
\end{quote}

\begin{quote}
\emph{āmcaṁ āṇi tyāñcaṁ javaḷcaṁ nātaṁ āhe.}
\end{quote}

\begin{quote}
We and they are closely related.
\end{quote}

Neuter. Plural: \mr{नाती}, \emph{nātī}.

Sentences write the neuter as \mr{नातं}. \mr{नातं} \mr{जोडणे}, to form a tie.

\begin{cousinweb}
\textbf{\mr{नाते}} (\emph{nāte}) is inherited: CDIAL lists Marathi \emph{nātẽ}, \textquotedblleft{}kinship, kinsman\textquotedblright{}, under the reconstructed \emph{*jñātra}, \textquotedblleft{}kinship\textquotedblright{} (CDIAL 5279).
\end{cousinweb}

\begin{grammarlens}[title={What you've built}]
One more word for describing what you lived and felt.
\end{grammarlens}

\subsection*{Your turn}
\begin{itemize}
\raggedright
  \item \emph{Say it:} \emph{nāte}
  \item \emph{Say it:} \emph{nāte}, once more
  \item \emph{Say it:} \emph{āmcaṁ āṇi tyāñcaṁ javaḷcaṁ nātaṁ āhe.} — We and they are closely related.
  \item \emph{Recall:} say the Marathi for to shiver with cold, then the Marathi for to limp, then say \emph{nāte} again
\end{itemize}

\begin{tcolorbox}[breakable,colback=teal!4,colframe=teal!35!black,title={Before you move on}]
What does \mr{नाते} mean? (\textquotedblleft{}A relationship; kinship\textquotedblright{}.) Say it once more.
\end{tcolorbox}
\section[\texorpdfstring{\mr{नातेवाईक} (nātevāīk)}{नातेवाईक (nātevāīk)}]{\mr{नातेवाईक} (nātevāīk) — a relative}

\label{lesson:MR-C494-natevaik}



\begin{quote}
Before the new one: say the Marathi for to limp, then the Marathi for a relationship.
\end{quote}

\subsection*{You'll want to know: \mr{नातेवाईक}}
\textbf{\mr{नातेवाईक}} — \emph{nātevāīk} — \textquotedblleft{}a relative\textquotedblright{}.

\begin{quote}
\mr{सुट्टीत} \mr{सगळे} \mr{नातेवाईक} \mr{आजीच्या} \mr{घरी} \mr{जमले}.
\end{quote}

\begin{quote}
\emph{suṭṭīt sagḷe nātevāīk ājīcyā gharī jamle.}
\end{quote}

\begin{quote}
In the holidays all the relatives gathered at Grandmother's house.
\end{quote}

Masculine.

Plural \mr{नातेवाईक}. Built on \mr{नाते}, kinship.

\begin{cousinweb}
\textbf{\mr{नातेवाईक}} (\emph{nātevāīk}) is a Marathi derivation of \mr{नाते}, \textquotedblleft{}kinship\textquotedblright{} (CDIAL 5279); the origin of the ending \emph{-vāīk} is not settled.
\end{cousinweb}

\begin{grammarlens}[title={What you've built}]
One more word for describing what you lived and felt.
\end{grammarlens}

\subsection*{Your turn}
\begin{itemize}
\raggedright
  \item \emph{Say it:} \emph{nātevāīk}
  \item \emph{Say it:} \emph{nātevāīk}, once more
  \item \emph{Say it:} \emph{suṭṭīt sagḷe nātevāīk ājīcyā gharī jamle.} — In the holidays all the relatives gathered at Grandmother's house.
  \item \emph{Recall:} say the Marathi for to limp, then the Marathi for a relationship, then say \emph{nātevāīk} again
\end{itemize}

\begin{tcolorbox}[breakable,colback=teal!4,colframe=teal!35!black,title={Before you move on}]
What does \mr{नातेवाईक} mean? (\textquotedblleft{}A relative\textquotedblright{}.) Say it once more.
\end{tcolorbox}
\section[\texorpdfstring{\mr{परिवार} (parivār)}{परिवार (parivār)}]{\mr{परिवार} (parivār) — a household; one's family}

\label{lesson:MR-C494-parivar}



\begin{quote}
Before the new one: say the Marathi for a relationship, then the Marathi for a relative.
\end{quote}

\subsection*{You'll want to know: \mr{परिवार}}
\textbf{\mr{परिवार}} — \emph{parivār} — \textquotedblleft{}a household; one's family\textquotedblright{}.

\begin{quote}
\mr{आमचा} \mr{परिवार} \mr{मोठा} \mr{आहे}.
\end{quote}

\begin{quote}
\emph{āmcā parivār moṭhā āhe.}
\end{quote}

\begin{quote}
Ours is a big family.
\end{quote}

Masculine.

More formal than \mr{कुटुंब}.

\begin{cousinweb}
\textbf{\mr{परिवार}} (\emph{parivār}) is a tatsama, taken from Sanskrit \emph{parivāra}, \textquotedblleft{}surroundings; followers\textquotedblright{} (CDIAL 7876).
\end{cousinweb}

\begin{grammarlens}[title={What you've built}]
One more word for describing what you lived and felt.
\end{grammarlens}

\subsection*{Your turn}
\begin{itemize}
\raggedright
  \item \emph{Say it:} \emph{parivār}
  \item \emph{Say it:} \emph{parivār}, once more
  \item \emph{Say it:} \emph{āmcā parivār moṭhā āhe.} — Ours is a big family.
  \item \emph{Recall:} say the Marathi for a relationship, then the Marathi for a relative, then say \emph{parivār} again
\end{itemize}

\begin{tcolorbox}[breakable,colback=teal!4,colframe=teal!35!black,title={Before you move on}]
What does \mr{परिवार} mean? (\textquotedblleft{}A household; one's family\textquotedblright{}.) Say it once more.
\end{tcolorbox}
\section[\texorpdfstring{\mr{अनाथ} (anāth)}{अनाथ (anāth)}]{\mr{अनाथ} (anāth) — an orphan}

\label{lesson:MR-C494-anath}



\begin{quote}
Before the new one: say the Marathi for a relative, then the Marathi for a household.
\end{quote}

\subsection*{You'll want to know: \mr{अनाथ}}
\textbf{\mr{अनाथ}} — \emph{anāth} — \textquotedblleft{}an orphan\textquotedblright{}.

\begin{quote}
\mr{त्या} \mr{संस्थेत} \mr{अनाथ} \mr{मुलांची} \mr{काळजी} \mr{घेतली} \mr{जाते}.
\end{quote}

\begin{quote}
\emph{tyā saṁsthet anāth mulāñcī kāḷjī ghetlī jāte.}
\end{quote}

\begin{quote}
Orphaned children are cared for at that institution.
\end{quote}

Masculine.

Also an adjective: \mr{अनाथ} \mr{मुलगी}. \mr{अनाथालय}, an orphanage.

\begin{cousinweb}
\textbf{\mr{अनाथ}} (\emph{anāth}) is a tatsama, taken from Sanskrit \emph{anātha}, \textquotedblleft{}having no master or protector\textquotedblright{} (Monier-Williams).
\end{cousinweb}

\begin{grammarlens}[title={What you've built}]
One more word for describing what you lived and felt.
\end{grammarlens}

\subsection*{Your turn}
\begin{itemize}
\raggedright
  \item \emph{Say it:} \emph{anāth}
  \item \emph{Say it:} \emph{anāth}, once more
  \item \emph{Say it:} \emph{tyā saṁsthet anāth mulāñcī kāḷjī ghetlī jāte.} — Orphaned children are cared for at that institution.
  \item \emph{Recall:} say the Marathi for a relative, then the Marathi for a household, then say \emph{anāth} again
\end{itemize}

\begin{tcolorbox}[breakable,colback=teal!4,colframe=teal!35!black,title={Before you move on}]
What does \mr{अनाथ} mean? (\textquotedblleft{}An orphan\textquotedblright{}.) Say it once more.
\end{tcolorbox}
\section[\texorpdfstring{\mr{सख्खा} (sakkhā)}{सख्खा (sakkhā)}]{\mr{सख्खा} (sakkhā) — full-blood (of a sibling)}

\label{lesson:MR-C494-sakkha}



\begin{quote}
Before the new one: say the Marathi for a household, then the Marathi for an orphan.
\end{quote}

\subsection*{You'll want to know: \mr{सख्खा}}
\textbf{\mr{सख्खा}} — \emph{sakkhā} — \textquotedblleft{}full-blood (of a sibling)\textquotedblright{}.

\begin{quote}
\mr{अरुण} \mr{माझा} \mr{सख्खा} \mr{भाऊ} \mr{आहे}.
\end{quote}

\begin{quote}
\emph{aruṇ mājhā sakkhā bhāū āhe.}
\end{quote}

\begin{quote}
Arun is my own brother.
\end{quote}

Agrees: \mr{सख्खा} \mr{भाऊ}, \mr{सख्खी} \mr{बहीण}. Also close: \mr{सख्खे} \mr{मित्र}.

\begin{cousinweb}
\textbf{\mr{सख्खा}} (\emph{sakkhā}) is a Marathi adjective whose origin is not settled.
\end{cousinweb}

\begin{grammarlens}[title={What you've built}]
One more word for describing what you lived and felt.
\end{grammarlens}

\subsection*{Your turn}
\begin{itemize}
\raggedright
  \item \emph{Say it:} \emph{sakkhā}
  \item \emph{Say it:} \emph{sakkhā}, once more
  \item \emph{Say it:} \emph{aruṇ mājhā sakkhā bhāū āhe.} — Arun is my own brother.
  \item \emph{Recall:} say the Marathi for a household, then the Marathi for an orphan, then say \emph{sakkhā} again
\end{itemize}

\begin{tcolorbox}[breakable,colback=teal!4,colframe=teal!35!black,title={Before you move on}]
What does \mr{सख्खा} mean? (\textquotedblleft{}Full-blood (of a sibling)\textquotedblright{}.) Say it once more.
\end{tcolorbox}
